%% slide11_study_note.tex
%% Detailed Study Note --- Slide 11: Ch 7 Adaptive Distance Protection (Zone Compensation)
%% TSA Meeting Preparation | Anoop V Eluvathingal | IIT Madras | 2026
\documentclass[12pt,a4paper]{article}

\usepackage[margin=2.5cm]{geometry}
\usepackage{amsmath,amssymb,bm}
\usepackage{booktabs,tabularx,array,longtable}
\usepackage{graphicx}
\usepackage{xcolor}
\usepackage{tcolorbox}
\usepackage{enumitem}
\usepackage{fancyhdr}
\usepackage{titlesec}
\usepackage{hyperref}
\usepackage{parskip}
\usepackage{algorithm}
\usepackage{algpseudocode}
\algrenewcommand\algorithmicrequire{\textbf{Input:}}
\algrenewcommand\algorithmicensure{\textbf{Output:}}
\newcommand{\kibr}{k_\mathrm{ibr}}
\newcommand{\Zapp}{Z_\mathrm{app}}
\newcommand{\Zfault}{Z_\mathrm{fault}}
\newcommand{\DZibr}{\Delta Z_\mathrm{ibr}}
\newcommand{\Zreach}{Z_\mathrm{reach}}

\definecolor{iitmnavy}{RGB}{15,40,80}
\definecolor{iitmgold}{RGB}{180,140,0}
\definecolor{lightblue}{RGB}{220,235,250}
\definecolor{lightyellow}{RGB}{255,252,220}
\definecolor{lightred}{RGB}{255,235,235}
\definecolor{lightgreen}{RGB}{230,255,230}
\definecolor{lightpurple}{RGB}{240,230,255}
\definecolor{lightorange}{RGB}{255,243,225}

\tcbuselibrary{skins,breakable}
\newtcolorbox{keybox}[1]{colback=lightblue,colframe=iitmnavy,
  fonttitle=\bfseries\color{white},title=#1,arc=3pt,boxrule=1pt,breakable}
\newtcolorbox{formulabox}[1]{colback=lightyellow,colframe=iitmgold,
  fonttitle=\bfseries,title=#1,arc=3pt,boxrule=1pt,breakable}
\newtcolorbox{resultbox}[1]{colback=lightgreen,colframe=green!50!black,
  fonttitle=\bfseries,title=#1,arc=3pt,boxrule=1pt,breakable}
\newtcolorbox{warnbox}[1]{colback=lightred,colframe=red!60!black,
  fonttitle=\bfseries,title=#1,arc=3pt,boxrule=1pt,breakable}
\newtcolorbox{archbox}[1]{colback=lightpurple,colframe=iitmnavy!60,
  fonttitle=\bfseries,title=#1,arc=3pt,boxrule=1pt,breakable}
\newtcolorbox{speakbox}{colback=orange!8,colframe=orange!70!black,
  fonttitle=\bfseries,title={What to Say at the TSA Meeting},arc=3pt,
  boxrule=1pt,breakable}

\pagestyle{fancy}
\fancyhf{}
\fancyhead[L]{\color{iitmnavy}\small\textbf{TSA Meeting Study Note}}
\fancyhead[R]{\color{iitmnavy}\small Slide 11 --- Ch~5 Adaptive Distance Protection}
\fancyfoot[C]{\thepage}
\fancyfoot[R]{\small\color{gray}Anoop V Eluvathingal $\cdot$ IIT Madras $\cdot$ 2026}
\renewcommand{\headrulewidth}{1.5pt}
\renewcommand{\headrule}{\color{iitmnavy}\hrule width\headwidth height\headrulewidth}

\titleformat{\section}{\large\bfseries\color{iitmnavy}}{}{0em}{}[\color{iitmnavy}\titlerule]
\titleformat{\subsection}{\normalsize\bfseries\color{iitmnavy!80}}{}{0em}{}

\hypersetup{colorlinks=true,linkcolor=iitmnavy,urlcolor=iitmnavy}

\begin{document}

%% ── Title Block ──────────────────────────────────────────────────────────────
\begin{tcolorbox}[colback=iitmnavy,coltext=white,arc=4pt,boxrule=0pt]
\begin{center}
  {\Large\bfseries TSA Meeting --- Slide 11 Study Note}\\[4pt]
  {\large Ch~5: Adaptive Distance Protection --- Zone Compensation for IBR}\\[6pt]
  {\normalsize Contributions C5 $\cdot$ Chapter 5 $\cdot$ SAMBP Framework}\\[2pt]
  {\small Anoop V Eluvathingal $\cdot$ IIT Madras $\cdot$ PhD Thesis 2026}
\end{center}
\end{tcolorbox}

\vspace{6pt}

%% ── 1. Slide Overview ────────────────────────────────────────────────────────
\section{Slide Overview and Narrative}

\begin{keybox}{What This Slide Shows}
Slide~11 presents \textbf{Ch~5: Adaptive Distance Protection}, covering Contributions C5 and C6. The slide's narrative has three parts:

\textbf{The problem:} Distance relays measure apparent impedance $\Zapp$. IBR reactive current injection during a fault ride-through adds a complex shift $\DZibr(\kibr, \varphi_\mathrm{ctrl})$ to the measured impedance. For lagging reactive injection (the dominant Volt-VAr control mode during voltage dips), the shift pushes $\Zapp$ \emph{into} Zone~1 for faults beyond 80\% of the line --- causing \textbf{overreach} (a false Zone~1 trip). For leading injection, the shift pushes $\Zapp$ \emph{out} of Zone~1 --- causing \textbf{underreach} (a missed trip).

\textbf{The solution:} The SAMBP-21 relay implements two complementary mechanisms: (1) an \textbf{adaptive quadrilateral Zone~1 boundary} that corrects for $\DZibr$ using the LM-estimated $\hat{\kibr}$, updated every 10~ms; and (2) \textbf{Zone~2 POTT acceleration} that bypasses the 300~ms Zone~2 timer when $\kibr > 0.5$, clearing faults within 32~ms via GOOSE permissive signalling.

\textbf{The result:} Overreach events fall from 17/50 to 2/50 (from 34\% to 4\%), underreach from 5/50 to 0/50, correct Zone~1 trips rise from 74\% to 96\%, with only 1~ms additional average trip time.
\end{keybox}

\subsection{Slide Key Results}

\begin{center}
\begin{tabular}{@{}lrrr@{}}
\toprule
\textbf{Metric} & \textbf{Fixed Relay} & \textbf{SAMBP-21} & \textbf{Change} \\
\midrule
Overreach events   & 17/50 (34\%) & 2/50  (4\%) & $-88\%$ relative \\
Underreach events  & 5/50  (10\%) & 0/50  (0\%) & eliminated \\
Correct Zone~1 trip & 74\%        & 96\%        & $+22$~pp \\
Avg trip time       & 18~ms       & 19~ms       & $+1$~ms only \\
\bottomrule
\end{tabular}
\end{center}

\textbf{Key observation from slide:} Zone adaptation nearly eliminates both overreach and underreach simultaneously, with negligible speed penalty. This is the central claim of Contribution~C5.

%% ── 2. Engineering Challenge ─────────────────────────────────────────────────
\section{Engineering Challenge: Why Fixed Zone~1 Fails for IBR}

\begin{warnbox}{The Core Problem: Apparent Impedance Distortion}
A distance relay at End~A of a transmission line measures:
\[
  \Zapp = \frac{\hat{V}_A}{\hat{I}_A}
\]
In a classical synchronous network, for a fault at fractional position $\alpha$ along the line with fault impedance $Z_F$:
\[
  \Zapp = \alpha Z_L + Z_F \left(1 + \frac{\hat{I}_B}{\hat{I}_A}\right)
\]
where $\hat{I}_B$ is the infeed current from the remote end. This is the classical infeed equation --- the relay always sees a \emph{larger} apparent impedance than the fault distance alone, never a smaller one.

\textbf{Why IBR breaks this:} A Type-4 IBR under Volt-VAr control injects reactive current $\hat{I}_B = \kibr I_r e^{j\beta}$ during the fault. The commanded angle $\beta$ depends on the control mode:
\begin{itemize}[nosep]
  \item \textbf{Volt-VAr absorb} (most common during nearby faults: $V_\mathrm{PCC}$ dips):
        $\beta = -90^\circ$ (lagging) $\Rightarrow$ $\Zapp$ shifts \emph{downward} in the $X$ direction
        $\Rightarrow$ faults at $\alpha > 0.8$ enter Zone~1 $\Rightarrow$ \textbf{overreach}.
  \item \textbf{Volt-VAr inject} (remote faults where $V_\mathrm{PCC}$ rises):
        $\beta = +90^\circ$ (leading) $\Rightarrow$ $\Zapp$ shifts \emph{upward}
        $\Rightarrow$ Zone~1 boundary missed $\Rightarrow$ \textbf{underreach}.
\end{itemize}
The conventional relay's Zone~1 boundary is set at 80--85\% of line length and is \emph{fixed}. It has no knowledge of the current IBR operating point and cannot adapt.
\end{warnbox}

\subsection{Magnitude of the Problem}

For a 100~km, 110~kV line with $Z_L = 0.35 + j0.40$~$\Omega$/km and a fault at $\alpha = 0.85$ (just outside Zone~1 in a classical network):
\[
  |\DZibr| = Z_F \kibr \frac{I_r}{I_A} = 10\,\Omega \times 0.5 \times 1.0 = 5\,\Omega
\]
The Zone~1 reactive reach is $X_\mathrm{reach} \approx 0.80 \times 0.40 \times 100 = 32\,\Omega$. A shift of $-5\,\Omega$ in $X$ moves the apparent impedance from $X_\mathrm{app} = 34\,\Omega$ (outside Zone~1) to $X_\mathrm{app} = 29\,\Omega$ (inside Zone~1) --- a clear overreach. This is not a corner case; at $\kibr = 0.5$ and $\beta = -90^\circ$, \textbf{every fault at $\alpha > 0.8$ overreaches}.

%% ── 3. Apparent Impedance Shift: Full Derivation ─────────────────────────────
\section{IBR Apparent Impedance Shift: Full Derivation}

\begin{formulabox}{Complete Apparent Impedance Equation}
The apparent impedance with IBR reactive current injection decomposes into three terms:
\begin{equation}
  \Zapp = \alpha Z_L \;+\; Z_F \;+\; \underbrace{Z_F \kibr \frac{I_r}{I_A} e^{j\beta}}_{\DZibr(\kibr,\,\beta)}
  \label{eq:zapp_full}
\end{equation}
The three terms represent:
\begin{enumerate}[nosep]
  \item $\alpha Z_L$: the line impedance to the fault (the correct measurement).
  \item $Z_F$: the fault resistance seen by End~A (present in classical networks too).
  \item $\DZibr = Z_F \kibr (I_r/I_A) e^{j\beta}$: the IBR infeed contribution (new in IBR networks).
\end{enumerate}

The reactive component of the shift is:
\begin{equation}
  \Delta X_\mathrm{ibr} = Z_F \kibr \frac{I_r}{I_A} \sin\beta
  \label{eq:dx_ibr}
\end{equation}

\textbf{Slide notation:} The slide writes $Z_\mathrm{app} = Z_\mathrm{fault} + \DZibr(\kibr, \varphi_\mathrm{ctrl})$ where $\varphi_\mathrm{ctrl}$ corresponds to $\beta$ (the commanded reactive angle). The ``radially outward shift'' mentioned on the slide refers specifically to the $\beta = -90^\circ$ case, where $\DZibr$ is a downward vector in the $R$--$X$ plane, moving the locus toward the origin (into Zone~1).
\end{formulabox}

\subsection{Why the Shift is Radially Outward from the Relay Perspective}

The relay at End~A is looking \emph{into} the line (toward increasing $\alpha Z_L$). ``Radially outward'' in the relay's impedance plane means toward larger $|Z|$ --- which corresponds to a fault appearing \emph{farther} away. However, for the lagging case ($\beta = -90^\circ$), $\Delta X < 0$, which means $\Zapp$ moves toward smaller $|Z|$ --- i.e., the fault appears \emph{closer} than it is. This is the \textbf{overreach} scenario. The phrase ``radially outward'' in the slide caption refers to the direction of the IBR shift vector in the complex plane (radially from the fault impedance point), not toward a larger apparent distance.

%% ── 4. Quadrilateral Characteristic ─────────────────────────────────────────
\section{Quadrilateral Distance Characteristic (Contribution C5)}

\begin{archbox}{Why Quadrilateral Instead of Mho}
The classical Mho (admittance circle) characteristic has two critical limitations:
\begin{enumerate}[nosep]
  \item \textbf{Zero resistive reach at origin}: The Mho circle passes through the origin. For high-resistance faults ($R_F > 0$), the fault locus runs along the $R$-axis --- outside the Mho circle. The Mho misses resistive faults entirely.
  \item \textbf{Reactive reach not independently controllable}: Extending the circle to cover resistive faults also extends the $X$ reach, increasing overreach risk.
\end{enumerate}
The \textbf{quadrilateral characteristic} provides independent control:
\begin{align}
  C_X &: \quad X_1 \leq X_\mathrm{relay} \leq X_\mathrm{reach} \label{eq:quad_x}\\
  C_R &: \quad R_1 \leq R_\mathrm{relay} \leq R_\mathrm{reach} \label{eq:quad_r}\\
  C_\mathrm{load} &: \quad R_\mathrm{relay} \geq \tan(\angle Z_L - \delta_\mathrm{load}) \cdot X_\mathrm{relay}
  \label{eq:quad_load}
\end{align}
The relay trips when all three conditions are simultaneously met: $T_{21} = \mathbf{1}[C_X \cap C_R \cap C_\mathrm{load}]$.

The load blinder $C_\mathrm{load}$ (set at $\delta_\mathrm{load} = 20^\circ$) ensures the operating zone does not overlap with the load impedance region (which also lies in the first quadrant of the $R$--$X$ plane at high $|Z|$).
\end{archbox}

\subsection{IBR Infeed Amplification and Resistance Reach Setting}

With IBR infeed, the apparent \emph{resistance} is amplified:
\begin{equation}
  R_\mathrm{app} = R_F \left(1 + \frac{1}{\kibr}\right)
  \label{eq:rfwd_amplification}
\end{equation}
For $\kibr = 0.1$ and $R_F = 0.004$~pu (arcing fault): $R_\mathrm{app} = 0.004 \times 11 = 0.044$~pu. The quadrilateral forward resistance reach must therefore be set to:
\[
  R_\mathrm{fwd,1} \geq R_F^\mathrm{max} \left(1 + \frac{1}{\kibr^\mathrm{min}}\right)
\]
For $\kibr^\mathrm{min} = 0.10$: $R_\mathrm{fwd,1} = 0.030$~pu covers all arc faults with $R_F \leq 0.004$~pu at any IBR penetration level $\kibr \geq 0.10$ (TR-34).

\textbf{Contrast with Mho}: At the Zone~1 circle boundary, the Mho has maximum resistance reach of approximately $R_\mathrm{mho} = X_\mathrm{reach} \cdot \tan(\angle Z_L)^{-1}$. For $\angle Z_L = 70^\circ$: $R_\mathrm{mho} \approx 0.36 X_\mathrm{reach}$. For arcing faults with high $R_F$ and IBR amplification, the Mho reach is exhausted before the quadrilateral.

%% ── 5. Adaptive Zone Compensation ───────────────────────────────────────────
\section{Adaptive Zone Compensation Mechanism}

\begin{formulabox}{SAMBP-21 Zone Boundary Update Law}
The Zone~1 reactive reach is adapted in real time:
\begin{equation}
  X_\mathrm{reach,adapt}(t) = X_\mathrm{reach,nom} - k_X \cdot \hat{\kibr}(t) \cdot \frac{I_\mathrm{rated}}{I_A}
  \label{eq:xreach_adapt}
\end{equation}
where:
\begin{itemize}[nosep]
  \item $X_\mathrm{reach,nom}$ = nominal Zone~1 reach (0.040~pu, set at 80\% of line reactance).
  \item $k_X = |Z_F^\mathrm{max}| \cdot |\sin\beta_\mathrm{max}| \approx 10\,\Omega \cdot 1 = 10\,\Omega$: reactive correction coefficient.
  \item $\hat{\kibr}(t)$ = real-time IBR penetration estimate from the Levenberg--Marquardt digital twin (Chapter~6).
  \item Update period: every \textbf{10~ms} (one power-frequency half-cycle).
\end{itemize}
The correction subtracts the expected reactive shift before it causes overreach, rather than detecting the overreach after the fact.
\end{formulabox}

\subsection{Zone~2 Acceleration Logic}

When the LM estimator detects $\hat{\kibr} > 0.5$, the IBR contribution is so dominant that Zone~1 instantaneous tripping is no longer reliable even with adaptation. SAMBP-21 then activates \textbf{Zone~2 POTT acceleration}:
\begin{itemize}
  \item The normal Zone~2 timer (200--300~ms) is bypassed.
  \item A permissive GOOSE signal is sent to End~B via IEC~61850.
  \item End~B confirms Zone~2 qualification and directional element \emph{forward}.
  \item Instantaneous trip is issued at End~B.
  \item End-to-end clearing: $4\,\mathrm{ms} + t_\mathrm{CT} + t_\mathrm{meas} \leq 8\,\mathrm{ms}$ (one cycle).
\end{itemize}
This is the \textbf{Permissive Overreaching Transfer Trip (POTT)} scheme. The 8~ms POTT trip time is well within the 100~ms memory polarisation window ($T_\mathrm{mem} = 100$~ms), ensuring directional integrity is maintained throughout.

\subsection{IBR-Secure POTT: Preventing Spurious Permissives}

A naive POTT can send a false permissive during a through-fault if IBR injection causes Zone~2 overreach. The IBR-secure POTT adds a $\kibr$ inhibit:
\begin{equation}
  T_\mathrm{POTT} = T_{Z2,A} \cap T_{Z2,B} \cap \mathbf{1}\!\left[\kibr < 0.45\right]
  \label{eq:pott_ibr}
\end{equation}
For $\kibr > 0.45$, the POTT reverts to Zone~2 with reach reduced to 110\% and timer set to 100~ms --- sacrificing some speed but maintaining security against through-fault overreach.

%% ── 6. Impedance Measurement and Cross-Polarisation ─────────────────────────
\section{Cross-Polarised Mho Element for Directional Security}

The SAMBP-21 uses a cross-polarised mho element for the directional supervisory function. The polarising voltage is derived from a healthy phase to eliminate memory-voltage dependency:
\begin{equation}
  \hat{V}_\mathrm{pol} = \hat{V}_b - \hat{V}_c \quad (\text{for phase~A fault})
  \label{eq:crosspol}
\end{equation}
The operate condition:
\begin{equation}
  \mathrm{Re}\!\left[\hat{I}_a \overline{Z_\mathrm{reach}} \cdot \overline{\hat{V}_\mathrm{pol}}\right] \geq 0
  \label{eq:crosspol_operate}
\end{equation}
Because $\hat{V}_{bc}$ remains near rated during single-phase and double-phase faults, the directional element does not lose polarisation as the memory voltage decays. This eliminates failure mode F6 (directional loss after memory expiry, which occurs at $t > T_\mathrm{mem}$ for classical self-polarised Mho).

\textbf{Timing margin:} Zone~1 trip time is $t_{Z1} \approx 20$~ms. Memory window is $T_\mathrm{mem} = 100$~ms. Margin: $T_\mathrm{mem}/t_{Z1} = 5\times$ --- the directional element remains valid even if the Mho element is delayed five-fold (e.g., due to CT saturation or communication delays).

%% ── 7. Results Explanation ───────────────────────────────────────────────────
\section{Results: What the Numbers Mean}

\begin{resultbox}{Interpreting the 50-Fault Simulation Results}
The 50-fault simulation covers the full ($\alpha$, $\kibr$, $\beta$) operating space:
\begin{itemize}[nosep]
  \item $\alpha \in \{0.1, 0.3, 0.5, 0.7, 0.85, 0.95\}$ (fault location)
  \item $\kibr \in \{0.1, 0.2, 0.3, 0.4, 0.5\}$ (IBR penetration)
  \item $\beta \in \{-90^\circ, 0^\circ, +90^\circ\}$ (control mode)
\end{itemize}

\medskip
\begin{tabular}{@{}lrrr@{}}
\toprule
\textbf{Metric} & \textbf{Fixed Relay} & \textbf{SAMBP-21} & \textbf{Improvement} \\
\midrule
Overreach events   & 17 (34\%) & 2 (4\%)  & $-88\%$ \\
Underreach events  & 5  (10\%) & 0 (0\%)  & $-100\%$ \\
Correct Zone~1 trip & 28 (56\%) & 48 (96\%) & $+40$~pp \\
Zone~2 clearance   & 22 (44\%) & 0  (0\%)  & eliminated (POTT handles) \\
Avg Zone~1 trip (ms) & 18 & 19 & $+1$ ms only \\
\bottomrule
\end{tabular}

\medskip
\textbf{The 2 residual overreach cases:} Both occur at $\kibr = 0.5$, $\alpha = 0.95$, $\beta = -90^\circ$ --- the extreme corner of the operating space where the IBR shift exceeds the adaptive correction's update bandwidth (10~ms update vs 8~ms fault transient). These cases are cleared by POTT within 32~ms.

\textbf{The $+1$~ms speed penalty:} The LM estimator requires 2--3 samples (10--15~ms) to converge to $\hat{\kibr}$. During this estimation period, Zone~1 asserts only after the first adaptive boundary update. This is the source of the 1~ms additional delay.
\end{resultbox}

\subsection{Coverage Map: TR-35 Figure~2}

The full coordination study (TR-35) maps the ($\kibr$, $\alpha$) plane into four protection regions:
\begin{center}
\begin{tabular}{@{}clll@{}}
\toprule
\textbf{Region} & \textbf{Condition} & \textbf{Primary function} & \textbf{Time} \\
\midrule
A & $\kibr < 0.06$ & 87L (line differential) only & 8--20~ms \\
B & $0.06 \leq \kibr$ and $\alpha \leq 0.80$ & Zone~2/POTT + 87L & 32~ms \\
C & $\alpha \leq 0.80$, $\kibr \leq 0.12$ & Zone~1 + TDCS & 20~ms \\
D & Remainder covered & Zone~1 + 87L & 20~ms \\
\bottomrule
\end{tabular}
\end{center}
\textbf{Key result from TR-35:} Zero coverage gaps across all 48 sampled $(\alpha, \kibr)$ cells. Every point in the operating space is covered by at least one protection function.

%% ── 8. Best Figure Recommendation ───────────────────────────────────────────
\section{Best Figure to Show at the TSA Meeting}

\begin{keybox}{Recommended Figure: TR-35 Figure~2 --- Full Protection Coverage Map}
\textbf{What it shows:} The $(\kibr, \alpha)$ plane divided into colour-coded regions (A, B, C, D) showing which combination of protection functions covers each operating point. The 87L channel (green band, $\kibr \geq 0.06$) runs in parallel with all distance elements and provides the universal back-stop.

\textbf{Why it is the best figure:}
\begin{enumerate}[nosep]
  \item It is the \emph{synthesis} result of Chapters 7--11: no single protection element covers the full plane; the coordinated set does.
  \item It shows \emph{zero gaps} --- a clear, definitive result.
  \item It directly answers ``what happens at high $\kibr$?'' and ``what happens for far-end faults?'' simultaneously.
  \item Examiners can read the regions directly without requiring detailed knowledge of any individual algorithm.
\end{enumerate}

\textbf{Where to find it:} TR-35/2026, page~3, Figure~2. Plot source: SAMBP test suite, coordination matrix module.

\textbf{Secondary figure:} TR-34 Figure~1 (R-X impedance plane with Quad Zone~1/Zone~2 rectangles, Mho circles, and IBR arc fault trajectory). Use this to explain \emph{why} the Mho misses resistive faults and why the quadrilateral is adopted. The IBR fault locus (red dashed line) running along the $R$-axis outside the Mho circle is immediately intuitive.

\textbf{Whiteboard supplement:} Draw the $R$--$X$ plane with: (1) a fixed Mho circle, (2) a fixed quadrilateral box, (3) an arrow showing the IBR lagging-injection shift vector, and (4) the \emph{adapted} quadrilateral box after SAMBP correction. This takes 60 seconds and makes the adaptation mechanism visually obvious.
\end{keybox}

%% ── 9. Cross-Thesis Connections ──────────────────────────────────────────────
\section{Cross-Thesis Connections}

\begin{archbox}{How This Chapter Connects to Other Work}
\textbf{Depends on (upstream):}
\begin{itemize}[nosep]
  \item \textbf{Ch~2 (IBR Characterisation)}: The failure modes F4 (overreach), F5 (underreach), F6 (directional loss) that this chapter solves are identified in Ch~2.
  \item \textbf{Ch~6 (Digital Twin / LM Estimator)}: The real-time $\hat{\kibr}$ estimate used in Eq.~\eqref{eq:xreach_adapt} comes from the EKF+LM digital twin. RMSE of $|\alpha - \hat{\alpha}| \leq 0.011$ (Ch~8) bounds the residual correction error.
  \item \textbf{Ch~8 (Cybersecurity)}: The GOOSE permissive signal in the POTT scheme is the attack surface addressed in Ch~9 (HMAC-SHA256 authentication, GOOSE replay protection).
\end{itemize}

\textbf{Enables (downstream):}
\begin{itemize}[nosep]
  \item \textbf{Ch~7 (WAPC)}: The POTT scheme's GOOSE infrastructure is reused by the Wide-Area Protection Coordinator for N-1 contingency coordination.
  \item \textbf{Ch~8 (Cybersecurity)}: The GOOSE authentication overhead (0.36~ms) is compatible with the POTT latency budget ($\leq 4$~ms GOOSE, 8~ms total).
  \item \textbf{TR-35}: This chapter's outputs (Zone~1/2 settings, POTT logic) feed directly into the full coordination study.
\end{itemize}

\textbf{SAMBP framework role:} This is Contribution~C5 (quadrilateral + IBR correction) and C6 (Monte Carlo reliability framework). Both feed into the SAMBP Stage-2 veto gate: if the adaptive reach correction cannot converge within one update cycle, Stage-2 blocks the Zone~1 trip and hands off to POTT (Zone~2). This is the SAMBP \emph{tiered response} --- attempt fast Zone~1, escalate to POTT, final backup via 87L.
\end{archbox}

%% ── 10. Key Equations Reference Card ────────────────────────────────────────
\section{Key Equations Reference Card}

\begin{formulabox}{Equations You Must Be Able to Write and Explain}

\textbf{1. Apparent impedance with IBR infeed (slide equation):}
\[
  \Zapp = \alpha Z_L + Z_F + Z_F \kibr \frac{I_r}{I_A} e^{j\beta}
\]
Say: ``$\DZibr$ is the IBR-induced shift. When $\beta = -90^\circ$ (lagging Volt-VAr), the shift has a negative imaginary component, pulling $\Zapp$ into Zone~1 for faults beyond the 80\% reach boundary.''

\textbf{2. Reactive shift:}
\[
  \Delta X = Z_F \kibr \frac{I_r}{I_A} \sin\beta
\]
Say: ``For $\beta = -90^\circ$, $\Delta X < 0$ --- the relay sees the fault as closer than it is. This is overreach.''

\textbf{3. Adaptive reach correction (SAMBP-21):}
\[
  X_\mathrm{reach,adapt}(t) = X_\mathrm{reach,nom} - k_X \hat{\kibr}(t) \frac{I_\mathrm{rated}}{I_A}
\]
Say: ``$k_X \approx 10\,\Omega$ is the maximum reactive correction coefficient. We subtract the expected shift before it causes a false trip.''

\textbf{4. Quadrilateral trip condition:}
\[
  T_{21} = \mathbf{1}\bigl[C_X \cap C_R \cap C_\mathrm{load}\bigr]
\]
Say: ``Independent $X$ and $R$ reach boundaries, plus a load blinder. The Mho cannot do this.''

\textbf{5. Resistive reach with IBR infeed amplification:}
\[
  R_\mathrm{app} = R_F \!\left(1 + \frac{1}{\kibr}\right)
\]
Say: ``At $\kibr = 0.1$, a 0.004~pu arc resistance appears as 0.044~pu to the relay. The quadrilateral forward reach is set to cover this.''

\textbf{6. IBR-secure POTT logic:}
\[
  T_\mathrm{POTT} = T_{Z2,A} \cap T_{Z2,B} \cap \mathbf{1}[\kibr < 0.45]
\]
Say: ``For $\kibr > 0.45$, POTT is inhibited and Zone~2 reverts to a 100~ms timer with 110\% reach to prevent through-fault false permissives.''

\textbf{7. POTT end-to-end trip time:}
\[
  t_\mathrm{total} = t_\mathrm{GOOSE} + t_\mathrm{CT} + t_\mathrm{meas} \leq 4 + 2 + 2 = 8\,\mathrm{ms}
\]
Say: ``32~ms on the slide includes a conservative GOOSE latency and relay processing time. The IEC~61850 P1 class guarantees $\leq 4$~ms GOOSE transfer.''
\end{formulabox}

%% ── 11. Speaker Script ───────────────────────────────────────────────────────
\section{Speaker Script (Timed to 3 Minutes)}

\begin{speakbox}
\textbf{[0:00--0:30] Open with the problem:}\\
``Distance protection is the backbone of transmission line protection worldwide --- it measures impedance rather than current, making it inherently selective. But IBR completely changes the measurement. When a Type-4 inverter injects reactive current during fault ride-through, it adds a complex impedance shift $\DZibr$ to whatever the relay measures. For lagging reactive injection --- which is exactly what Volt-VAr control does during a nearby voltage dip --- this shift pulls the apparent impedance into Zone~1 for faults beyond the 80\% reach boundary. The relay declares an instantaneous trip for a fault that is not in its zone. That is overreach.''

\textbf{[0:30--1:15] Explain the solution:}\\
``SAMBP-21 addresses this with two mechanisms. First, the adaptive boundary: instead of fixing Zone~1 at 80\% of the line, we update the reactive reach boundary every 10~ms based on the real-time $\hat{\kibr}$ estimate from the LM digital twin. The correction is $X_\mathrm{adapt} = X_\mathrm{nom} - k_X \hat{\kibr}$. We pre-compensate for the shift before it causes a false trip. Second, when IBR penetration is high --- above 50\% --- Zone~1 adaptation alone is insufficient, because the estimator convergence time approaches the fault transient duration. In that regime, SAMBP-21 immediately asserts a Zone~2 POTT permissive over IEC~61850 GOOSE to the remote end, achieving clearance in under 32~ms --- versus the conventional 200--300~ms Zone~2 timer.''

\textbf{[1:15--2:00] Present the results:}\\
``On 50 simulated faults spanning the full IBR operating envelope, the fixed relay had 17 overreach events and 5 underreach events --- 44\% of faults were misclassified. SAMBP-21 reduces overreach to 2 events and underreach to zero. Correct Zone~1 clearance rises from 74\% to 96\%. The average trip time increases by only 1~ms, from 18 to 19~ms. Speed and selectivity improve simultaneously.''

\textbf{[2:00--2:30] Point to TR-35 coverage map:}\\
``The full coordination study in TR-35 maps the ($\kibr$, $\alpha$) operating space into four protection regions. The key result is that every point in the space --- every combination of IBR penetration and fault location --- is covered by at least one protection function. There are zero gaps. Below $\kibr = 0.06$, the 87L line differential is the primary element. Between 0.06 and 0.12, Zone~1 with TDCS. Above 0.12, Zone~2/POTT plus 87L. The coloured map summarises the entire distance protection coordination result in a single figure.''

\textbf{[2:30--3:00] Close with the contribution:}\\
``Contributions C5 and C6: the IBR-corrected quadrilateral eliminates overreach without sacrificing resistive fault coverage; the Zone~2 POTT with IBR inhibit provides sub-cycle clearance for high-penetration scenarios. Together, these restore distance protection reliability in IBR-dominated networks to the same standard expected in synchronous networks. The Monte Carlo reliability study in C6 quantifies this: detection probability 0.993 versus 0.917 for the classical Mho.''
\end{speakbox}

%% ── 12. Anticipated Q and A ──────────────────────────────────────────────────
\section{Anticipated Questions and Answers}

\begin{keybox}{Q1: Why not simply lower the Zone 1 reach setting to avoid overreach?}
\textbf{Answer:} Reducing the Zone~1 reach from 80\% to, say, 65\% would push the overreach risk boundary back, but at the cost of leaving 35\% of the line unprotected by fast instantaneous tripping. All faults in the ``gap'' (65--80\% of the line) would then require Zone~2 clearance (200--300~ms), which exceeds the generator ride-through capability and the IBR memory polarisation window of 100~ms. The adaptive approach maintains 80\% coverage when $\kibr$ is low and adjusts only when IBR penetration is high enough to cause a real overreach risk. It also inherently handles the underreach case (leading injection) by moving the boundary in the opposite direction.
\end{keybox}

\begin{keybox}{Q2: The POTT requires communication --- what if the GOOSE link fails?}
\textbf{Answer:} SAMBP-21 includes a GOOSE watchdog: if no heartbeat is received from the remote end within 2~ms (IEC~61850 P1 timeout), the relay reverts to the local Zone~2 timer (200~ms). This is the same fallback as a conventional distance relay --- the POTT provides acceleration, but the relay remains functional without it. The cybersecurity study (Ch~12/TR-52) shows that HMAC-SHA256 authentication adds only 0.36~ms overhead to each GOOSE frame, well within the 4~ms latency budget. GOOSE flooding attacks are mitigated by rate limiting at the switch layer.
\end{keybox}

\begin{keybox}{Q3: How does the relay know the IBR commanded angle beta?}
\textbf{Answer:} The relay does not need to know $\beta$ explicitly. The LM digital twin estimates $\hat{\kibr}$ from the measured terminal voltages and currents (the EKF+LM algorithm in Ch~8). The adaptive reach correction uses only $\hat{\kibr}$ and $I_\mathrm{rated}/I_A$ --- it does not require $\beta$. The correction coefficient $k_X$ is pre-set conservatively at the worst-case angle ($\beta = -90^\circ$): $k_X = |Z_F^\mathrm{max}| = 10\,\Omega$. If $\beta$ is less adverse (e.g., $\beta = 0^\circ$), the correction over-compensates slightly, making Zone~1 slightly conservative but not causing underreach (the residual error is $\Delta X_\mathrm{residual} \leq 0.055\,\Omega$, negligible).
\end{keybox}

\begin{keybox}{Q4: What about three-phase faults --- do they overreach too?}
\textbf{Answer:} Yes, but less severely. For a balanced three-phase fault, the IBR injects balanced reactive current, so the apparent impedance shift is purely reactive ($\Delta X$ term only, no $\Delta R$ component). The overreach condition is identical to the single-phase case: $\Delta X < 0$ for lagging injection. However, three-phase faults are rarer and typically produce larger differential currents (87L trips faster), so the Zone~1 overreach window is shorter. The 50-fault simulation includes three-phase faults; the adaptive correction handles them identically.
\end{keybox}

\begin{keybox}{Q5: How is Zone 2 POTT different from Distance Protection with pilot wires?}
\textbf{Answer:} Classic pilot-wire protection uses dedicated copper wires or a dedicated fibre channel for a differential current comparison signal. POTT uses the existing IEC~61850 substation LAN (GOOSE) --- no dedicated channel is needed. GOOSE is a multicast Ethernet frame with sub-4~ms guaranteed latency on a P1-class substation LAN. The POTT logic is entirely implemented in software within the distance relay, making it retrofit-compatible with existing IEC~61850 substations. The key difference from a pilot differential scheme is that POTT is a \emph{permission} signal (``End~A has seen a Zone~2 fault''), not a measurement signal --- it does not require time-synchronised current samples.
\end{keybox}

\vspace{1em}
\begin{center}
\begin{tcolorbox}[colback=iitmnavy!10,colframe=iitmnavy,arc=4pt,boxrule=1pt,width=0.85\textwidth]
\centering
\textbf{Chapter~5 One-Line Summary for Examiners}\\[4pt]
\textit{``IBR reactive injection distorts apparent impedance by $\DZibr = Z_F \kibr e^{j\beta}$. SAMBP-21 pre-compensates with a real-time adaptive Zone~1 boundary updated every 10~ms from the LM estimate of $\hat{\kibr}$, and accelerates clearance via POTT when $\kibr > 0.5$. Overreach falls from 34\% to 4\%; correct Zone~1 trips rise from 74\% to 96\%; speed penalty is 1~ms.  TR-35 shows zero coverage gaps across the full ($\alpha$, $\kibr$) operating space.''}
\end{tcolorbox}
\end{center}

\end{document}
